\documentclass[11pt,a4paper]{article}

% ======================================================================
% CONFIGURACIÓN DEL PROYECTO Y PAQUETES
% ======================================================================
\usepackage[utf8]{inputenc}
\usepackage[T1]{fontenc}
\usepackage[spanish,es-tabla]{babel}
\usepackage[margin=2.5cm, headheight=30pt]{geometry} 
\usepackage{graphicx}
\usepackage{fancyhdr}
\usepackage{xcolor}
\usepackage[colorlinks=true, linkcolor=blue, urlcolor=blue]{hyperref}
\usepackage{bookmark}

% --- Matemáticas y Electrónica ---
\usepackage{amsmath, amssymb}
\usepackage{siunitx}
\usepackage[american]{circuitikz}

% --- Elementos de Diseño y Código ---
\usepackage{tcolorbox}
\usepackage{booktabs}
\usepackage{enumitem}
\usepackage{listings}
\usepackage{tabularx}

% ======================================================================
% MACROS Y ESTILOS PERSONALIZADOS
% ======================================================================
\sisetup{
    output-decimal-marker = {,},
    per-mode = symbol,
    inter-unit-product = \ensuremath{{}\cdot{}}
}

\lstset{
    language=Python,
    basicstyle=\ttfamily\footnotesize,
    keywordstyle=\color{blue}\bfseries,
    stringstyle=\color{teal},
    commentstyle=\color{green!50!black}\itshape,
    backgroundcolor=\color{gray!5},
    frame=single,
    rulecolor=\color{gray!40},
    breaklines=true,
    showstringspaces=false,
    numbers=left,
    numberstyle=\tiny\color{gray},
    numbersep=5pt
}

\newtcolorbox{info}[1][]{
    colback=blue!5!white,
    colframe=blue!75!black,
    fonttitle=\bfseries,
    title=Información Importante,
    #1
}

\newtcolorbox{warning}[1][]{
    colback=red!5!white,
    colframe=red!75!black,
    fonttitle=\bfseries,
    title=Advertencia de Instrumentación,
    #1
}

\newcommand{\tecla}[1]{\colorbox{gray!20}{\texttt{\footnotesize\textbf{#1}}}}
\newcommand{\menu}[1]{\textbf{\textsf{#1}}}

% ======================================================================
% ESTILO DE PÁGINA
% ======================================================================
\pagestyle{fancy}
\fancyhf{}
\lhead{\textbf{Circuitos Electrónicos II}\\ \small Guía de Laboratorio N.º 1}
\rhead{Ingeniería Mecatrónica -- UCB Tarija\\ \small Semestre 1-2026}
\cfoot{Página \thepage}
\renewcommand{\headrulewidth}{0.4pt}

\begin{document}

% ======================================================================
% CABECERA Y METADATOS DEL DOCUMENTO
% ======================================================================
\begin{center}
    \rule{\textwidth}{1pt}\\[0.3cm]
    {\LARGE \textbf{GUÍA DE LABORATORIO N.º 1}}\\[0.2cm]
    {\Large \textbf{Caracterización Dinámica de Impedancia y Puente de Wheatstone en CA para un Sensor de Temperatura}}\\[0.2cm]
    \rule{\textwidth}{1pt}
\end{center}

\vspace{0.2cm}
\noindent\textbf{CARRERA:} Ingeniería Mecatrónica\\
\textbf{ASIGNATURA:} Circuitos Electrónicos II\\
\textbf{PROFESOR:} M.Sc. Ing. Bernardo Quiroga Turdera\\
\textbf{MODALIDAD:} Trabajo Colaborativo en Equipos de 3 Integrantes

\vspace{0.5cm}
% ======================================================================
% 1. RESULTADOS DE APRENDIZAJE
% ======================================================================
\section{RESULTADOS DE APRENDIZAJE (ABET / CDIO)}
Al finalizar esta práctica de laboratorio, el estudiante será capaz de:
\begin{itemize}[label=$\bullet$]
    \item \textbf{ABET Student Outcome 1:} Modelar matemáticamente la respuesta frecuencial y la impedancia compleja de un sensor de temperatura resistivo con parásitos de línea en régimen sinusoidal.
    \item \textbf{ABET Student Outcome 6:} Conducir experimentación en corriente alterna, adquirir datos con osciloscopio y procesarlos mediante scripts de Python para validar el comportamiento teórico.
    \item \textbf{CDIO (Design \& Implement):} Diseñar y balancear un Puente de Wheatstone alimentado en CA para eliminar tensiones de offset de CC en la instrumentación de temperatura del Proyecto Final Integrador (PFI).
\end{itemize}

% ======================================================================
% 2. CONTEXTO INDUSTRIAL
% ======================================================================
\section{CONTEXTO INDUSTRIAL Y ¿POR QUÉ USAMOS UN PUENTE DE WHEATSTONE?}
En plantas industriales (químicas, alimentarias, manufactura 4.0), las variables físicas como la temperatura se miden mediante sensores resistivos (como termistores NTC o PT100). Sin embargo, la variación de resistencia ($\Delta R$) ante un cambio de temperatura de \qty{1}{\celsius} es pequeña en comparación con la resistencia nominal del sensor ($R_0$).

\subsection*{El Problema del Divisor de Tensión Simple}
Si conectáramos el termistor en un divisor de tensión común alimentado con \qty{5}{\volt}, obtendríamos un voltaje de salida con un offset de CC gigantesco (por ejemplo, \qty{2.5}{\volt}) y una variación de señal útil microscópica (por ejemplo, \qty{10}{\milli\volt}). Si intentamos amplificar esa pequeña señal por $100$ con un amplificador operacional, el circuito intentará entregar $2.5\text{ V} \times 100 = \qty{250}{\volt}$, saturando e inutilizando el amplificador.

\subsection*{La Solución: El Puente de Wheatstone}
El Puente de Wheatstone utiliza dos divisores de tensión en paralelo para realizar una resta analógica de voltaje en origen.

\begin{center}
    \begin{circuitikz}[scale=1.0, transform shape]
        % Rediseñado para dar holgura a las etiquetas
        \draw (-1.5,4) to[V, l=$V_{CC} (\qty{5}{\volt})$] (-1.5,0);
        \draw (-1.5,4) -- (2.5,4) -- (6.5,4);
        \draw (-1.5,0) -- (2.5,0) -- (6.5,0);
        \draw (2.5,0) node[ground]{};
        
        \draw (2.5,4) to[R, l=$R_1$ (\qty{10}{\kilo\ohm})] (2.5,2) node[circle,fill,inner sep=1.5pt]{};
        \draw (2.5,2) to[R, l=$R_2$ (\qty{10}{\kilo\ohm})] (2.5,0);
        
        \draw (6.5,4) to[R, l=$R_3$ (\qty{10}{\kilo\ohm})] (6.5,2) node[circle,fill,inner sep=1.5pt]{};
        \draw (6.5,2) to[vR, l=$R_{NTC} (R_0 + \Delta R)$] (6.5,0);
        
        % Etiquetas de Nodos más claras
        \node[left, align=center] at (2.5,2) {\textbf{Nodo A}\\(Fijo)};
        \node[right, align=center] at (6.5,2) {\textbf{Nodo B}\\(Sensor)};
        
        \draw (2.5,2) to[short, -o] (4,2) node[above]{$+$};
        \draw (6.5,2) to[short, -o] (5,2) node[above]{$-$};
        \node at (4.5, 2) {$V_{diff}$};
    \end{circuitikz}
\end{center}

\begin{itemize}
    \item \textbf{Nodo A (Referencia fija):} Entrega \qty{2.50}{\volt}.
    \item \textbf{Nodo B (Sensor NTC):} Entrega \qty{2.50}{\volt} $\pm \Delta V$.
\end{itemize}
La salida diferencial entre ambos nodos ($V_{diff} = V_A - V_B$) cancela la ``manta'' continua de \qty{2.5}{\volt}. A la temperatura de reposo, el puente está balanceado ($V_{diff} = \qty{0}{\volt}$). Al variar la temperatura, el puente genera una tensión de salida diferencial centrada en cero de pocos milivoltios, lista para ser amplificada por un Amplificador de Instrumentación sin saturarlo.

\subsection*{¿Por qué lo alimentamos en Corriente Alterna (CA)?}
Al alimentar el puente con una señal senoidal en CA de alta frecuencia (\qty{10}{\kilo\hertz}), nos inmunizamos contra las derivas térmicas de corriente continua y permitimos el uso de filtrado activo frecuencial. Sin embargo, los cables de extensión largos presentan capacitancias parásitas ($C_{para}$) que introducen un desfase angular ($\theta$) que debe ser modelado matemáticamente.

% ======================================================================
% 3. MATERIALES
% ======================================================================
\section{MATERIALES E INSTRUMENTACIÓN REQUERIDA}

\noindent\textbf{Por Equipo de Trabajo (3 Integrantes):}
\begin{itemize}[label=$\circ$, itemsep=0pt]
    \item $1 \times$ Termistor NTC de \qty{10}{\kilo\ohm} a \qty{25}{\celsius} ($\beta = \qty{3950}{\kelvin}$).
    \item $3 \times$ Resistores de película metálica de \qty{10}{\kilo\ohm} al 1\% de tolerancia.
    \item $1 \times$ Resistor de película metálica de \qty{10}{\ohm} (simulación de resistencia de cable $R_{cable}$).
    \item $1 \times$ Inductor de \qty{15}{\micro\henry} (o sustituto pasivo disponible $L_{cable}$).
    \item $1 \times$ Capacitor de poliéster de \qty{4.7}{\nano\farad} (simulación de capacitancia parásita de línea $C_{para}$).
    \item $1 \times$ Protoboard, cables jumpers, termómetro digital o multímetro con sonda de temperatura.
\end{itemize}

\noindent\textbf{Equipamiento de Banco de Laboratorio:}
\begin{itemize}[label=$\circ$, itemsep=0pt]
    \item $1 \times$ Generador de funciones senoidales.
    \item $1 \times$ Osciloscopio digital con puerto USB para exportación de datos CSV.
    \item $1 \times$ Laptop con Python 3.10+ (\texttt{pandas}, \texttt{numpy}, \texttt{matplotlib}) y LTSpice instalado.
\end{itemize}

% ======================================================================
% 4. FUNDAMENTACIÓN MATEMÁTICA
% ======================================================================
\section{FUNDAMENTACIÓN MATEMÁTICA Y CÁLCULOS PREVIOS}
\begin{info}
Antes de armar el circuito, el equipo debe resolver las siguientes ecuaciones en su libreta de laboratorio. Estos cálculos son prerequisito para ingresar al laboratorio.
\end{info}

\noindent\textbf{Resistencia de la NTC (Steinhart-Hart reducida):}
\begin{equation}
    R_{NTC}(T) = R_0 \cdot e^{\beta \left( \frac{1}{T} - \frac{1}{T_0} \right)}
\end{equation}
\textit{Calcular $R_{NTC}$ para $T = 20\ ^\circ\text{C}, 30\ ^\circ\text{C}, 40\ ^\circ\text{C}, 50\ ^\circ\text{C}$ y $60\ ^\circ\text{C}$ ($T_0 = \qty{298.15}{\kelvin}$, $R_0 = \qty{10}{\kilo\ohm}$, $\beta = \qty{3950}{\kelvin}$).}

\vspace{0.3cm}
\noindent\textbf{Impedancia Compleja de la Rama del Sensor $Z_{sensor}(j\omega)$:}
\begin{equation}
    Z_{sensor}(j\omega) = R_{cable} + j\omega L_{cable} + \frac{R_{NTC}(T)}{1 + j\omega C_{para} R_{NTC}(T)}
\end{equation}
\textit{Para $f_c = \qty{10}{\kilo\hertz}$ ($\omega_c = 2\pi \cdot 10000$), expresar $Z_{sensor}$ en forma polar ($\vert{}Z\vert{} \angle \theta$) para $T = \qty{30}{\celsius}$.}

\vspace{0.3cm}
\noindent\textbf{Tensión de Desbalance Fasorial del Puente $\hat{V}_{out}$:}
\begin{equation}
    \hat{V}_{out}(j\omega) = \hat{V}_{in} \left( \frac{Z_{sensor}(j\omega)}{R_3 + Z_{sensor}(j\omega)} - \frac{R_2}{R_1 + R_2} \right)
\end{equation}
\textit{Para $\hat{V}_{in} = 1\text{ V}_{pico} \angle 0^\circ$ (\qty{2}{\volt}$_{pp}$), calcular la magnitud en voltios y la fase en grados de $\hat{V}_{out}$ a $T = \qty{30}{\celsius}$.}

% ======================================================================
% 5. PROCEDIMIENTO EXPERIMENTAL
% ======================================================================
\section{PROCEDIMIENTO EXPERIMENTAL (PASO A PASO)}

\begin{center}
    \begin{circuitikz}[scale=1.0, transform shape]
        % Se separó la fuente hacia la izquierda para evitar solapamientos con Nodo A
        \draw (-2,4) to[sV, l=$V_{in}$ (CA)] (-2,0);
        \draw (-2,4) -- (2.5,4) -- (6.5,4);
        \draw (-2,0) -- (2.5,0) -- (6.5,0);
        \draw (2.5,0) node[ground]{};
        
        \draw (2.5,4) to[R, l=$R_1$] (2.5,2) node[circle,fill,inner sep=1.5pt]{};
        \draw (2.5,2) to[R, l=$R_2$] (2.5,0);
        \node[left, align=center] at (2.5,2) {Nodo A\\(Canal 1)};
        
        \draw (6.5,4) to[R, l=$R_3$] (6.5,2) node[circle,fill,inner sep=1.5pt]{};
        \draw (6.5,2) to[generic, l=$Z_{sensor}$ (NTC + Parásitos)] (6.5,0);
        \node[right, align=center] at (6.5,2) {Nodo B\\(Canal 2)};
        
        \draw (2.5,2) to[short, -o] (4,2) node[above]{$+$};
        \draw (6.5,2) to[short, -o] (5,2) node[above]{$-$};
        \node at (4.5, 2) {$\hat{V}_{out}$};
    \end{circuitikz}
\end{center}

\subsection*{Paso 1: Verificación de Componentes (Modelo F.I.V.E.)}
Mida con el multímetro de precisión los valores reales de $R_1, R_2, R_3$ y $R_0$ (NTC a temperatura ambiente de laboratorio). Anote las desviaciones respecto a los valores nominales.

\subsection*{Paso 2: Ensamblaje en Protoboard}
\begin{enumerate}
    \item Monte el Puente de Wheatstone según el esquemático.
    \item En la rama 4 ($Z_{sensor}$), conecte el termistor NTC en \textbf{paralelo} con el capacitor parásito $C_{para} = \qty{4.7}{\nano\farad}$, y todo esto en \textbf{serie} con $R_{cable} = \qty{10}{\ohm}$ y $L_{cable} = \qty{15}{\micro\henry}$.
\end{enumerate}

\subsection*{Paso 3: Instrumentación y Configuración}
\begin{warning}
Nunca conecte la pinza de tierra (caimán negro) de las sondas del osciloscopio en un punto que no sea GND. Hacerlo provocará un cortocircuito.
\end{warning}
\begin{enumerate}
    \item Conecte el Generador de Funciones a la entrada del puente ($V_{in}$). Ajuste una onda senoidal de \qty{2}{\volt}$_{pp}$ (\qty{1}{\volt}$_{pico}$) a una frecuencia de \qty{10}{\kilo\hertz}.
    \item Conecte el \textbf{Canal 1 ($CH1$)} del osciloscopio en el Nodo A (Referencia).
    \item Conecte el \textbf{Canal 2 ($CH2$)} del osciloscopio en el Nodo B (Sensor).
    \item Utilice la función matemática del osciloscopio (\menu{MATH}: \tecla{CH1 - CH2}) para visualizar la señal diferencial limpia $\hat{V}_{out}$.
\end{enumerate}

\subsection*{Paso 4: Adquisición de Datos}
Para 5 puntos de temperatura diferentes (\qty{20}{\celsius}, \qty{30}{\celsius}, \qty{40}{\celsius}, \qty{50}{\celsius}, \qty{60}{\celsius}) calentando suavemente el termistor y monitoreando con el termómetro:
\begin{enumerate}
    \item Mida la \textbf{amplitud pico} de la señal de salida $V_{out}$ (traza MATH).
    \item Active los \textbf{cursores de tiempo} en el osciloscopio y mida el retraso temporal ($\Delta t$) entre el pico de $V_{in}$ (CH1) y el pico de $V_{out}$ (MATH).
    \item \textbf{Exportación:} Guarde la traza en la memoria USB en formato \texttt{.csv} para cada punto de temperatura.
\end{enumerate}

% ======================================================================
% 6. VALIDACIÓN EN PYTHON
% ======================================================================
\newpage
\section{VALIDACIÓN COMPUTACIONAL EN PYTHON}
Cada equipo deberá ejecutar el siguiente script en Python (\texttt{lab1\_validation.py}) para procesar los archivos \texttt{.csv} exportados del osciloscopio y calcular los errores porcentuales respecto al modelo teórico:

\begin{lstlisting}[language=Python]
import numpy as np
import pandas as pd
import matplotlib.pyplot as plt

# ==============================================================================
# 1. CARGA DE DATOS EXPERIMENTALES (CSV del Osciloscopio)
# ==============================================================================
# Archivo CSV esperado con columnas: 'Time', 'Channel1' (Vin), 'Channel2' (VB)
data = pd.read_csv("lab1_medicion_30C.csv")

t = data['Time'].values
v_in = data['Channel1'].values
v_b = data['Channel2'].values

# Tension diferencial medida: Vout = VA - VB (VA es Vin / 2)
v_a = v_in / 2.0
v_out_exp = v_a - v_b

# ==============================================================================
# 2. PROCESAMIENTO DE SENAL
# ==============================================================================
v_out_pico_exp = np.max(v_out_exp)

# Calculo del desfase temporal (Delta t) por deteccion de picos
idx_pico_in = np.argmax(v_in)
idx_pico_out = np.argmax(v_out_exp)
delta_t_exp = t[idx_pico_out] - t[idx_pico_in]

f_c = 10000.0  # Frecuencia de portadora 10 kHz
fase_grados_exp = delta_t_exp * f_c * 360.0

# ==============================================================================
# 3. MODELO TEORICO ANALITICO
# ==============================================================================
T_celsius = 30.0
T_kelvin = T_celsius + 273.15
R_0 = 10000.0; Beta = 3950.0; T_0 = 298.15

# Resistencia teorica NTC
R_ntc = R_0 * np.exp(Beta * (1.0/T_kelvin - 1.0/T_0))

# Impedancias
omega = 2 * np.pi * f_c
Z_c = 1.0 / (1j * omega * 4.7e-9)
Z_l = 1j * omega * 15e-6
R_cable = 10.0

Z_paralelo = (R_ntc * Z_c) / (R_ntc + Z_c)
Z_sensor = R_cable + Z_l + Z_paralelo

# Salida Fasorial Teorica
V_in_fasor = 1.0  # 1V pico
V_out_teorico = V_in_fasor * (Z_sensor / (10000.0 + Z_sensor) - 0.5)

v_out_pico_teorico = np.abs(V_out_teorico)
fase_grados_teorica = np.angle(V_out_teorico, deg=True)

# ==============================================================================
# 4. CALCULO DE ERRORES Y GRAFICACION
# ==============================================================================
err_amp = np.abs(v_out_pico_exp - v_out_pico_teorico) / v_out_pico_teorico * 100.0
err_fase = np.abs(fase_grados_exp - fase_grados_teorica) / np.abs(fase_grados_teorica) * 100.0

print(f"--- RESULTADOS DE VALIDACION (T = {T_celsius} °C) ---")
print(f"Amplitud Pico Teorica:  {v_out_pico_teorico:.4f} V | Experimental: {v_out_pico_exp:.4f} V | Error: {err_amp:.2f}%")
print(f"Fase Angular Teorica:   {fase_grados_teorica:.2f}° | Experimental: {fase_grados_exp:.2f}° | Error: {err_fase:.2f}%")

plt.figure(figsize=(10, 5))
plt.plot(t * 1e6, v_in, 'b--', label="V_in (Portadora 10kHz)", alpha=0.6)
plt.plot(t * 1e6, v_out_exp, 'r-', label="V_out Diferencial Experimental")
plt.xlabel("Tiempo (us)")
plt.ylabel("Tension (V)")
plt.title(f"Validacion de Salida Diferencial del Puente a {T_celsius} °C")
plt.grid(True)
plt.legend()
plt.savefig(f"validacion_puente_{int(T_celsius)}C.png")
plt.show()
\end{lstlisting}

% ======================================================================
% 7. RÚBRICA
% ======================================================================
\section{RÚBRICA ANALÍTICA DE EVALUACIÓN (100 PUNTOS)}
El informe técnico en formato IEEE debe entregarse de forma digital incluyendo el código de Python en un repositorio Git.

\begin{table}[h!]
    \centering
    \renewcommand{\arraystretch}{1.3} 
    \begin{tabularx}{\textwidth}{|>{\raggedright\arraybackslash}p{2.5cm}|>{\raggedright\arraybackslash}X|>{\raggedright\arraybackslash}X|>{\raggedright\arraybackslash}X|c|}
        \hline
        \textbf{Criterio Evaluativo} & \textbf{Sobresaliente (90-100 pts)} & \textbf{Aceptable (70-89 pts)} & \textbf{Insuficiente (0-69 pts)} & \textbf{Peso} \\
        \hline
        \textbf{Fundamentación Teórica} & Demostración matemática completa de $Z_{sensor}(j\omega)$ y $\hat{V}_{out}$. Explicación impecable de la cancelación de offset de CC por el puente. & Deducción analítica con errores menores de álgebra compleja o redondeo de constantes. & Copia de fórmulas sin desarrollo de nodos en CA; no comprende por qué se usa el puente. & \textbf{30\%} \\
        \hline
        \textbf{Simulación LTSpice} & Simulación con análisis \textit{AC Sweep} y \textit{Transient} incorporando componentes parásitos reales de línea. & Simulación correcta pero utilizando componentes pasivos ideales sin parásitos. & Simulación inoperativa o con errores de convergencia en SPICE. & \textbf{20\%} \\
        \hline
        \textbf{Implementación Física} & Armado pulcro; balance del puente verificado; uso correcto de la resta matemática ($CH1-CH2$) en el osciloscopio. & Armado funcional pero ruidoso; mediciones con escalas de osciloscopio de baja resolución. & Circuito inoperativo; incapacidad de medir el desfase o mala conexión de tierra ($GND$). & \textbf{20\%} \\
        \hline
        \textbf{Validación en Python} & Script automatizado que importa CSVs, calcula la fase por detección de picos y cuantifica los errores porcentuales. & Script de Python funcional pero que requiere ingresar los datos de pico manualmente en el código. & No presenta código de Python o grafica datos ingresados a mano en procesadores simples. & \textbf{20\%} \\
        \hline
        \textbf{Informe Técnico IEEE} & Formato IEEE impecable; análisis crítico de causas de error (tolerancias, EMI); código alojado en GitHub. & Formato IEEE con errores de diagramación o gráficas de osciloscopio sin rotulación de ejes. & Formato libre o descuidado; entrega tardía sin justificación. & \textbf{10\%} \\
        \hline
    \end{tabularx}
\end{table}

\end{document}